\section{Ordinary flags from the original local period data}
\label{sec:ordinary-flags-v44}

The purpose of this chapter is the full original implication
\(\mathrm{LiteratureRankNHypothesis}\Rightarrow\mathrm{OrdFlag}\).
Its conclusion is an actual ordered invariant flag over
\(\overline{\mathbb Q}_\ell\), together with potential crystallinity
for every original local bridge.  Equality of characteristic
polynomials alone is not substituted for that conclusion.

Fix a number field \(F\), its original algebraic closure \(\Omega\),
a place \(w\mid\ell\), a prolongation \(A\), and the original
field isomorphism \(\iota:\overline{\mathbb Q}_\ell\simeq\mathbb C\).
Write \(D\), \(I\), \(W\) for the decomposition, inertia, and Weil
groups.  The local Artin map is the source map
\(a:W\to F_w^\times\), with
\(\operatorname{ord}_w(a(s))=-\deg(s)\).
Every period assertion below is made using the original tied model
\(P=\mathrm{PP}(\ell)\) and an arbitrary inducing bridge
\(B_w\); finite coefficient extensions may be enlarged inside
\(\overline{\mathbb Q}_\ell\).  The isomorphism \(\iota\) is used
only algebraically and is never assumed continuous.

The retained published inputs are precisely the original global Artin
reciprocity and algebraic Hecke-character theorems
\texttt{N.aut.art}, \texttt{N.aut.chA}, \texttt{N.aut.chW}, the
original Chevalley unit-congruence input \texttt{N.loc.chev}, and
Berger's de Rham--potentially semistable theorem
\texttt{N.aut.berger}.  We also retain the original period-ring
foundations: finite coefficient models; comparison isomorphisms and
their tensor, dual, exactness, and finite-extension descent laws;
the fundamental identity
\[
 B_{\mathrm{cris}}^{\varphi=1}\cap
             \operatorname{Fil}^0B_{\mathrm{dR}}=\mathbb Q_\ell;
 \tag{OFL.1}
\]
and the crystalline comparison for an unramified rank-one character,
with Hodge--Tate weight zero and its actual Weil character.
These are the same foundations of the original \texttt{LocalPublishedTheorems},
\texttt{PstLaws}, and \texttt{TwistAbove} period models.  They are
listed as retained foundations, rather than counted as new proofs.
The intermediate invariant-vector, weight-order, and flag assertions
used by the original source are all proved below.

\subsection{Rank-one characters, with their exact local normalization}

For an integer-valued function \(h\) on embeddings of \(F\), put
\[
 A_h(x)=\prod_{\tau:F_w\hookrightarrow\overline{\mathbb Q}_\ell}
       \tau(x)^{h(\tau|_F)}\quad(x\in F_w^\times).
 \tag{OFL.2}
\]
There are finitely many continuous embeddings.  Consequently the
product is nonzero, is multiplicative, and satisfies
\(A_{h+k}=A_hA_k\) and \(A_{-h}=A_h^{-1}\).
The literal \(\mathrm{OrdChar}(\iota,a,\omega,h,\psi)\)
says that \(\omega:F_w^\times\to\mathbb C^\times\) is smooth,
\(\psi:D\to\overline{\mathbb Q}_\ell^\times\) is continuous, and
\[
       \psi(s)=\iota^{-1}\bigl(\omega(a(s))\bigr)A_{-h}(a(s))
                         \quad(s\in W).
 \tag{OFL.3}
\]

% BEGIN OFL1
\paragraph{OFL1. The inertia algebraic part and inversion.}
\label{claim:OFL1}
An ordinary character is locally algebraic on inertia, with weight
\(h\).  Its inverse is an ordinary character with data
\((\omega^{-1},-h,\psi^{-1})\).
\begin{proof}
On inertia, multiply (OFL.3) by the character
\(\iota^{-1}(\omega\circ a)^{-1}\).  This factor has open
kernel: its kernel is the inverse image of \(\ker\omega\)
under the continuous map \(a|_I\).  Thus the product is exactly
the algebraic character \(A_{-h}\), which is the original
\(\mathrm{IsWtPart}\) convention.  If another allowed Artin map
is used, their original \(\mathrm{glob}\) clauses identify the
images of the local ideles after global Artin.  We verify the
needed injectivity in this profinite target.  Suppose the local
idele of \(x\in F_w^\times\) lies in the Artin kernel.  The
original global reciprocity law identifies that kernel with the
closure of the diagonal \(F^\times\) times the connected
infinity subgroup.  The latter has finite components one.
Consequently, for every positive integer \(a\), we can approximate
\(x\) arbitrarily closely at \(w\) by \(k\in F^\times\)
such that at every other finite place both \(k\) and \(k^{-1}\)
are integral and \((k-1)/a\) is integral.  Indeed these are open
conditions in the restricted-product topology, also imposed on
the inverse idele, and the given local idele satisfies them.

First \(x\) must be a local unit.  If \(|x|_w<1\), sufficiently
close approximants have a fixed positive \(w\)-valuation \(d\)
and valuation zero elsewhere.  Their principal ideal is
\(w^d\), so their integral norm has fixed absolute value
\(R=q^d\), where \(q=\operatorname{Norm}(w)\).
Choose \(a=q(R+1)+1\).  It is prime to \(w\) and greater
than \(R+1\).  A close approximant is integral everywhere,
and \((k-1)/a\) is integral also at \(w\), because \(a\)
is a local unit there.  Thus \(k=1+ay\) with \(y\in\mathcal O_F\).
The determinant defining its norm gives
\(\operatorname{Norm}_{F/\mathbb Q}(k)\equiv1\pmod a\).
An integer of absolute value \(R\) with this congruence must
equal one, contradicting \(d>0\).  Applying the same argument
to \(x^{-1}\) excludes \(|x|_w>1\).

Now fix \(N\ge1\), and let \(m\ge1\) annihilate the finite
unit group \((\mathcal O_F/q^N\mathcal O_F)^\times\).
The original \texttt{N.loc.chev} gives a positive integer
\(a\), prime to \(q\), such that every global unit congruent
to one modulo \(a\) is an \(m\)-th power of a global unit.
Choose an approximant with
\(|k-x|_w\le |q|_w^N<1\).  It is a unit at every finite
place, and \((k-1)/a\) is integral everywhere, by the same
prime-to-\(w\) argument.  Hence \(k=u^m\) for a global unit
\(u\), so \(k\equiv1\pmod{q^N}\).  The ultrametric
inequality yields \(|x-1|_w\le |q|_w^N\).
As \(N\) is arbitrary and \(|q|_w<1\), we obtain \(x=1\).
This proves injectivity after global Artin itself.  Applying it
to the quotient of the two local Artin values makes the two maps
equal, and supplies the universal Artin-map quantifier in
\(\mathrm{IsWtPart}\).  Taking inverses in (OFL.3)
and using (OFL.2) proves the second assertion.  The inverse has the
same smooth kernel and a continuous coefficient-field character.
\end{proof}
% END OFL1

% BEGIN OFL2
\paragraph{OFL2. A CM field realizing the fixed local field.}
\label{claim:OFL2}
For a finite extension \(K/\mathbb Q_\ell\) in the original
coefficient closure and the embedding \(j:\Omega\hookrightarrow
\overline{\mathbb Q}_\ell\) of a bridge, there is a CM subfield
\(E\subset\Omega\) whose image generates \(K\) over
\(\mathbb Q_\ell\).  It can be chosen with
\(0\ne\delta\in E\) such that every complex embedding sends
\(\delta\) to a purely imaginary number and \(j(\delta)\in
\mathbb Q_\ell\).
\begin{proof}
Choose a primitive element \(\alpha\) of \(K\).  The separability
of its minimum polynomial gives an open \(\ell\)-adic neighborhood
of its coefficients in which a root remains sufficiently close to
\(\alpha\) to generate \(K\).  To check the last assertion,
take disjoint balls around its conjugates; Hensel's lemma supplies
the nearby root, and Krasner's argument says that an automorphism
fixing that root must fix the nearest conjugate \(\alpha\).
Equality of the degrees then gives equality of the generated fields.
Independently choose a real monic polynomial of the same degree
having distinct real roots.  Having that many distinct real roots is
an open real condition, by disjoint intervals and sign changes.
Weak approximation for rational coefficients supplies one monic
polynomial satisfying both open conditions.  Its nearby root
\(\beta\in K\) is algebraic over \(\mathbb Q\), all its
conjugates are real, and \(\mathbb Q_\ell(\beta)=K\).
Lift \(\beta\) to \(b\in\Omega\) through \(j\): the
rational polynomial splits in \(\Omega\), so its image has every
root there.

The polynomial \(X^2-X+\ell\) has a simple root near zero in
\(\mathbb Z_\ell\), since its derivative there is \(-1\).
If that root is \(z\), then \(d=2z-1\) satisfies
\(d^2=1-4\ell\).  Choose a square root \(\delta\) of
\(1-4\ell\) in \(\Omega\), changing its sign so that
\(j(\delta)=d\).  Since \(1-4\ell<0\), every complex image
of \(\delta\) is nonzero and purely imaginary.
Now \(E=\mathbb Q(b,\delta)\) is the quadratic totally imaginary
extension of the totally real field \(\mathbb Q(b)\).  It is CM,
and its local image generates exactly \(K\).  This construction
also works at \(\ell=2\); no exclusion of the original prime
\(3\) occurs.
\end{proof}
% END OFL2

% BEGIN OFL3
\paragraph{OFL3. Globalizing the prescribed local algebraic weight.}
\label{claim:OFL3}
For every bridge \(B_w\) and every weight \(h\), a global
algebraic Hecke character on a field as in OFL2 supplies a continuous
local Galois character \(\nu\) with that same locally algebraic
weight.  The bridge can retain exactly its local field, embedding
\(j\), and coefficient splitting field.
\begin{proof}
Use OFL2 with the local field of \(B_w\).  Restrict the valuation
to \(E\) to obtain a place \(u\); the completion \(E_u\)
and \(F_w\) identify with the same finite field \(K\) inside
the coefficient closure.  Restriction through these identifications
identifies their decomposition groups.  Thus changing the global
field does not change the local field or period model.

Let \(S\) be the embeddings \(E\to\overline{\mathbb Q}_\ell\)
induced by embeddings of \(K\).  Every such embedding sends
\(\delta\) to the same nonzero \(d\in\mathbb Q_\ell\).
After composition with \(\iota\), complex conjugation changes
that image to its negative.  Hence \(\iota S\) and its conjugate
are disjoint.  An embedding in \(S\) determines a unique embedding
of \(F_w\): agreement on the generating field \(E\) determines
the map on \(K\).  Assign infinity exponent \(-h(\tau|_F)\)
on \(\iota S\), the opposite exponent on its conjugate, and
zero on the remaining conjugate pairs.  This gives an integral
infinity type \(n\) with \(n(\phi)+n(\bar\phi)=0\).

Here is the usual existence argument for a Hecke character of this
type, so that globalization is not a new local supplier assumption.
On units of a CM field, the algebraic expression prescribed by
\(n\) has absolute value one under every conjugate: complex
conjugation pairs the exponents with their negatives.  It is an
algebraic unit, so its values are roots of unity.  They lie in one
fixed number field and therefore form a finite group.  Choose a
finite congruence level on which all these unit values are one;
this can be seen by embedding that finite group in a local field
and choosing a sufficiently small neighborhood of one.  The
infinity character and the trivial character on this finite-level
compact subgroup then agree on their intersection with diagonal
\(E^\times\).  They define a character on the subgroup they
generate together with \(E^\times\), trivial on \(E^\times\).
Extend it to the idele group using divisibility of \(\mathbb C^\times\).
Its restriction to the stated open subgroup is continuous; hence
the extension is continuous.  This is an algebraic Hecke character
with exactly the prescribed infinity type.  No conductor restriction
is needed for the present endpoint.

The retained algebraic Hecke-character theorem supplies its
\(\ell\)-adic Galois character \(\mu\).  Pull it to \(D\)
through the common local field to define \(\nu\).  The Artin
compatibility in that theorem gives inertia weight \(h\), since
the chosen exponent was \(-h\).  To check compatibility of the
two global fields, take their compositum in \(\Omega\): its
chosen completion is still \(K\), and both restriction maps
induce the identity there.  Local reciprocity on the compositum
therefore identifies the two inertia formulas.  The induced
embeddings have the same period multiplicities because the
underlying local representation and local field are literally the
same.  This proves all assertions, including the original bridge
quantifiers.
\end{proof}
% END OFL3

% BEGIN OFL4
\paragraph{OFL4. The full rank-one \texttt{DRat} and potential-crystalline law.}
\label{claim:OFL4}
If \(\mathrm{OrdChar}(\iota,a,\omega,h,\psi)\), then
\[
 \mathrm{DRat}\!\left(\iota,w,\operatorname{rk}_1\psi,
                 \tau\mapsto\{h(\tau)\},
                 s\mapsto X-\omega(a(s))\right),
 \qquad \mathrm{PotCrys}(\operatorname{rk}_1\psi).
 \tag{OFL.4}
\]
These assertions hold for every original inducing bridge.
\begin{proof}
Fix such a bridge.  OFL1 and OFL3 give a globalized local character
\(\nu\) with the same inertia algebraic weight as \(\psi\).
Their quotient \(\eta=\psi\nu^{-1}\) has open inertia kernel:
intersect the two open kernels of their smooth inertia parts.
There is a finite extension of the local field whose inertia is
in that kernel.  Indeed an open identity neighborhood in inertia
contains the intersection with an open normal subgroup of the
absolute Galois group; the latter corresponds to a finite extension.
On that extension \(\eta\) is unramified.  The original
rank-one crystalline comparison makes it de Rham with jump zero.
Finite-extension descent preserves being de Rham and its labeled
jump multiplicities.  The retained Hecke-character theorem
\texttt{chA}, applied to \(\nu\), gives jump \(h(\tau)\).
Tensoring with \(\eta\) gives the same jumps for \(\psi\).
Induced embeddings of \(F\) extend uniquely to \(F_w\), as
in UPA11, so this proves the exact \(\mathrm{IsJumpsAt}\)
quantifier, not only a list indexed by a chosen embedding family.

The retained algebraic-character WD theorem \texttt{chW}, together
with the unramified rank-one comparison and its finite-extension
descent, identifies the parameter of \(\psi\) with the smooth
character
\[
 \xi(s)=\iota\bigl(\psi(s)A_h(a(s))\bigr).
 \tag{OFL.5}
\]
More explicitly, this identity holds first for each globalized
weight character: it is the Artin formula in \texttt{chW} on
the common local field.  Tensor products and inverse characters
give arbitrary integral weights.  For the quotient \(\eta\)
the Weil factor is its actual character \(\iota\eta|_W\),
since the period comparison for a potentially unramified character
uses the original Weil action.  Multiplying these identities yields
(OFL.5); finite descent recovers all inertia elements and not just
a Frobenius power.  This also verifies the source degree sign.
By (OFL.3), (OFL.5) is exactly \(\omega(a(s))\), so its
characteristic polynomial is the polynomial in (OFL.4).

Finally apply \texttt{N.aut.berger} to this de Rham representation.
Its WD space has dimension one, and a nilpotent endomorphism of a
one-dimensional space is zero.  The retained equivalence between
zero monodromy and potential crystallinity therefore proves the
second assertion.  The bridge was arbitrary throughout.
\end{proof}
% END OFL4

\subsection{Distinct inertia and a fixed period}

% BEGIN OFL5
\paragraph{OFL5. The finite-inertia character projectors.}
\label{claim:OFL5}
Let \(I\to\Gamma\) be normal in the sense that its image is
preserved by conjugation.  Let \(r\) be a complex representation
with finite inertia image, and suppose
\[
 \det(X-r(g))=\prod_{j=0}^{n-1}(X-\mu_j(g))
                                 \quad(g\in\Gamma),
 \tag{OFL.6}
\]
where the characters \(\mu_j|_I\) are pairwise distinct.
Then inertia is the direct sum of its one-dimensional
\(\mu_j\)-spaces.  Every nilpotent endomorphism commuting with
inertia is zero.
\begin{proof}
If \(r(s)=1\) for \(s\in I\), (OFL.6) gives
\(\mu_j(s)=1\) for all \(j\).  Thus all these characters
factor through the finite group \(H=r(I)\).  For their induced
characters \(\chi_j\), form
\[
        e_j=|H|^{-1}\sum_{h\in H}\chi_j(h)^{-1}r(h).
 \tag{OFL.7}
\]
Here \(r(h)\) denotes the faithful action of the image group.
The finite character sum of a nontrivial character is zero:
translate the sum by an element on which that character is not
one.  It follows that \(e_j e_k=\delta_{jk}e_j\) and that
the image of \(e_j\) is the \(\chi_j\)-eigenspace.
The trace identity obtained from (OFL.6) shows
\(\operatorname{tr}(e_j)=1\).  Since an idempotent has
eigenvalues zero and one, its image has dimension one.
The degree in (OFL.6) gives \(\dim r=n\); hence these \(n\)
lines sum to the entire space.  A commuting nilpotent preserves
each line, acts there by a nilpotent scalar, and is zero on all
lines.  This proves the assertion, including the empty rank-zero
case.
\end{proof}
% END OFL5

% BEGIN OFL6
\paragraph{OFL6. The inertia eigenline has its prescribed full Weil character.}
\label{claim:OFL6}
Under OFL5, \(r(g)\) acts on the \(j\)-th inertia line by
\(\mu_j(g)\), for every \(g\in\Gamma\).
\begin{proof}
Normality of the inertia image and the fact that \(\mu_j\)
is a character of all of \(\Gamma\) show that its inertia
eigenspace is preserved by every \(g\).  Let its scalar there
be \(c_j(g)\).  Taking the trace after applying (OFL.7) gives
\[
 c_j(g)=\operatorname{tr}(r(g)e_j)
  =|H|^{-1}\sum_{h\in H}\chi_j(h)^{-1}
                        \sum_k\mu_k(g)\chi_k(h)
  =\mu_j(g).
\]
Consequently every vector fixed by inertia is fixed by the whole
Weil group when one of the \(\mu_j\) is the trivial character:
that character has the only trivial inertia restriction, by
distinctness.  In the WD application, normality is checked directly
from \(\deg(gsg^{-1})=0\).
\end{proof}
% END OFL6

% BEGIN OFL7
\paragraph{OFL7. Descending the invariant vector through WD scalarization.}
\label{claim:OFL7}
In the original potentially semistable period model, use the raw
complex Weil action \(\mathrm{wdRho}\) on \(\mathrm{WDSpace}\),
before Frobenius semisimplification.  A nonzero
inertia-fixed vector of this complex space yields a nonzero
inertia-fixed vector of \(D_{\mathrm{pst}}\).  If all the former
inertia-fixed vectors are Weil fixed, so are all the latter.
\begin{proof}
The original scalarization map \(\mathrm{wdMk}:D_{\mathrm{pst}}
\to\mathrm{WDSpace}\) is additive, intertwines the period Weil
action with this raw \(\mathrm{wdRho}\) action,
and has complex spanning image.  The inertia action on the period
space has open kernel, hence finite image.  Average over this finite
image.  If that average were zero on every period vector, its
complexification would vanish on the spanning image and therefore
on the whole WD space.  On a nonzero inertia-fixed WD vector,
however, the unnormalized average is multiplication by the positive
integer equal to the group order.  This contradiction gives a
nonzero inertia-fixed period vector.

For the second assertion, let \(x\) be any such period vector.
The period Frobenius \(\varphi\) commutes with the corrected
Weil action.  Hence \(\varphi^m x\) is inertia fixed for each
\(m\ge0\), and its scalarized image is fixed by every Weil
element.  For \(z=\mathrm{weilD}(s)x-x\) we obtain
\[
             \mathrm{wdMk}(\varphi^m z)=0\quad(m\ge0).
 \tag{OFL.8}
\]
The original faithfulness of all these Frobenius translates implies
\(z=0\).  To recall that faithfulness, scalarization selects one
embedding component of the maximal unramified coefficient algebra;
successive Frobenius maps visit every component over any finite
unramified subextension containing the coordinates of \(z\).
Vanishing on all components is vanishing of \(z\).  Thus one
does not incorrectly assume that a single scalarization map is
injective.  This proves the assertion on the original period space.
\end{proof}
% END OFL7

% BEGIN OFL8
\paragraph{OFL8. Turning a Weil-fixed period into a Frobenius-fixed period.}
\label{claim:OFL8}
Suppose \(D_{\mathrm{pst}}\) contains a nonzero inertia-fixed
vector and every inertia-fixed vector is Weil fixed.  Then it
contains \(z\ne0\) fixed both by \(\varphi\) and by the
entire local Galois action.
\begin{proof}
Let \(x\ne0\) be inertia fixed.  Every period vector is fixed
by an open local Galois subgroup, so some positive power of inverse
Frobenius fixes \(x\) under the Galois action.  It also fixes
\(x\) under the corrected Weil action.  Since that corrected
action is the Galois action times
\(\varphi^{-\alpha(s)}\), with
\(\alpha(s)=f_w\deg(s)\), it follows that
\(\varphi^m x=x\) for some integer \(m\ge2\).
Replacing the initial exponent by twice itself ensures this bound.

Choose \(c\) in the maximal unramified extension of
\(\mathbb Q_\ell\) with Frobenius orbit of length exactly
\(m\).  For example take a primitive
\((\ell^m-1)\)-st root of unity.  Frobenius sends it to
its \(\ell\)-th power, and
\(\ell^i<\ell^m-1\) for \(0\le i<m\), \(m\ge2\),
so the indicated orbit has that length.  For integers \(j\ge0\),
set
\[
                  z_j=\sum_{i=0}^{m-1}\varphi^i(c^j x).
 \tag{OFL.9}
\]
The identity \(\varphi^m(c)=c\) and
\(\varphi^m x=x\) makes \(\varphi z_j=z_j\).
At least one of \(z_0,\ldots,z_{m-1}\) is nonzero: otherwise
the Vandermonde matrix in the distinct scalars
\(c,\varphi(c),\ldots,\varphi^{m-1}(c)\) annihilates the
vector \((x,\varphi x,\ldots,\varphi^{m-1}x)\), forcing
\(x=0\).  This argument works in the vector space over the
maximal unramified field and does not require a choice of a
Frobenius eigenvalue in \(\mathbb C\).

Each \(z_j\) is inertia fixed: inertia fixes the unramified
scalars, commutes with \(\varphi\), and fixes \(x\).
The chosen nonzero \(z_j\) is therefore Weil fixed as well.
Since it is already Frobenius fixed, the corrected and uncorrected
Weil actions agree on it.  Its stabilizer for the original Galois
action is open and closed, and the Weil group is dense in the
decomposition group.  It is thus fixed by every Galois element.
\end{proof}
% END OFL8

% BEGIN OFL9
\paragraph{OFL9. Nonnegative Hodge--Tate weights recover a Galois invariant.}
\label{claim:OFL9}
Let \(V\) be a de Rham representation over a finite coefficient
field, with all its labeled Hodge--Tate weights nonnegative in the
source convention.  If \(0\ne z\in D_{\mathrm{pst}}(V)\)
is killed by monodromy and fixed by both Frobenius and Galois, then
\(V\) has a nonzero Galois-fixed vector.
\begin{proof}
Galois invariance puts \(z\) in \(D_{\mathrm{st}}(V)\).
The kernel of monodromy there is the image of
\(D_{\mathrm{cris}}(V)\), so write its unique antecedent as
\(y\).  The crystalline-to-semistable inclusion is injective and
commutes with Frobenius and Galois.  Thus \(y\) is fixed by both.
Its de Rham image lies in \(D_{\mathrm{dR}}(V)\).

Here is the filtration step with its sign.  Enlarge coefficients
to contain every embedding of the local field.  The source labeled
multiplicity is the dimension of the corresponding graded
filtration piece.  Every negative graded piece is zero by the
nonnegative-weight assumption.  A finite-dimensional, separated,
exhaustive decreasing filtration therefore satisfies
\(D_{\mathrm{dR}}=\operatorname{Fil}^0D_{\mathrm{dR}}\):
start at a sufficiently negative index where the filtration is the
whole space and use equality of successive pieces up to zero.
Descent through the enlarged coefficient field preserves this
identity.

Choose a \(\mathbb Q_\ell\)-basis of \(V\), and write
\(y=\sum_i v_i\otimes b_i\).  Its Frobenius invariance says
\(\varphi(b_i)=b_i\), and its membership in filtration zero
says \(b_i\in\operatorname{Fil}^0B_{\mathrm{dR}}\).
The fundamental identity (OFL.1) gives \(b_i\in\mathbb Q_\ell\).
Hence \(y=v\otimes1\) for a unique \(v\in V\).
The nonzero and invariant assertions descend through this injective
tensor map.  This gives \(v\ne0\) and \(g(v)=v\) for every
local Galois element, as required.
\end{proof}
% END OFL9

% BEGIN OFL10
\paragraph{OFL10. The complete local core of the slope argument.}
\label{claim:OFL10}
Let \(r:D\to\mathrm{GL}_n(\overline{\mathbb Q}_\ell)\)
be continuous, and assume the original \(\mathrm{DRat}\) with
weights \(H(\tau)\) and polynomials
\(\prod_j(X-\mu_j(s))\), where the smooth Weil characters
\(\mu_j\) have pairwise distinct inertia restrictions.  Then
\(r\) is potentially crystalline for every original bridge.
If some \(\mu_j=1\) and all \(H(\tau)\) are nonnegative,
there is a nonzero vector fixed by every element of \(D\).
\begin{proof}
For an arbitrary bridge the \(\mathrm{DRat}\) datum gives
de Rhamness and the asserted WD polynomials.  Berger's theorem
gives a potentially semistable finite coefficient model.  We
apply OFL5--OFL8 to its raw complex Weil action
\(\mathrm{wdRho}\), which is the action actually intertwined
by \(\mathrm{wdMk}\).  The semisimplified action
\(\mathrm{wdRep}\) defining \(\mathrm{wdPst}\) differs from
it by a commuting unipotent factor raised to \(\deg(s)\).
Commuting matrices are simultaneously triangularizable over
\(\mathbb C\); this factor has diagonal entries one, so the
two actions have the same characteristic polynomial on every
Weil element.  On inertia the factor is the identity.  Thus the
\(\mathrm{DRat}\) polynomials in the hypothesis are exactly
the raw-action polynomials required in OFL5--OFL6, and its
inertia image is finite.  No invariant vector is being inferred
merely from a Frobenius-semisimplified action.
The monodromy operator is
nilpotent and commutes with inertia, since inertia has degree zero
in the WD commutation identity.  OFL5 therefore makes that
operator zero.  The original monodromy-zero comparison gives
potential crystallinity, including the specified finite model.

When \(\mu_j=1\), OFL5--OFL6 give the nonzero trivial
inertia line and trivial Weil action on it.  OFL7--OFL8 produce
the nonzero period fixed by Frobenius and Galois.  OFL9 gives a
nonzero invariant vector in the finite coefficient model.
Embedding its coordinates in \(\overline{\mathbb Q}_\ell\)
preserves nonzeroness and the matrix identities.  This yields
the invariant of \(r\) itself.  For this existence assertion
one may choose any inducing bridge, whose existence is part of
the original local model; potential crystallinity was proved for
each bridge separately.  Rank zero is harmless: the latter
assertion is then vacuous because there is no index \(j\).
\end{proof}
% END OFL10

\subsection{Twisting, quotienting, and the ordered flag}

% BEGIN OFL11
\paragraph{OFL11. The exact \texttt{DRat} twist formula.}
\label{claim:OFL11}
Suppose \(r\) has \(\mathrm{DRat}\) weights \(H(\tau)\)
and polynomials \(\prod_j(X-\alpha_j(s))\), and a continuous
rank-one character \(\eta\) has weights \(c(\tau)\) and
polynomials \(X-\beta(s)\).  Then
\(r'(g)=\eta(g)r(g)\) is continuous and has weights
\(\{h+c(\tau):h\in H(\tau)\}\) and polynomials
\[
                    \prod_j(X-\alpha_j(s)\beta(s)).
 \tag{OFL.10}
\]
\begin{proof}
Continuity follows entry by entry, and the inverse matrix is
\(\eta(g)^{-1}r(g)^{-1}\), so it is continuity into the
original general linear group.  Work over a common finite
coefficient extension containing the two models.  Tensoring a de
Rham representation with a de Rham line preserves its comparison
isomorphism.  A filtered line with jump \(c\) shifts every
filtration index by \(c\), giving exactly the displayed multiset.
Berger makes the representations potentially semistable.  On
their period modules tensoring with the line multiplies the Weil
matrix by its scalar \(\beta(s)\); the monodromy on a line
is zero.  A scalar multiple of a matrix with roots
\(\alpha_j(s)\) has roots \(\alpha_j(s)\beta(s)\), with
multiplicities, so its monic characteristic polynomial is
(OFL.10).  These identifications commute with all inducing
embeddings and with coefficient enlargement.  They therefore
prove the full original \(\mathrm{DRat}\), for every bridge.
\end{proof}
% END OFL11

% BEGIN OFL12
\paragraph{OFL12. Removing a character line preserves the remaining data.}
\label{claim:OFL12}
Suppose \(r\) contains an invariant line with character
\(\psi_0\).  Assume that \(r\) has \(\mathrm{DRat}\)
weights \(\{h_0(\tau),\ldots,h_{n-1}(\tau)\}\) and
polynomials \(\prod_{j=0}^{n-1}(X-\alpha_j(s))\), and the
line has weights \(h_0\) and polynomial \(X-\alpha_0(s)\).
Then the quotient by this line is continuous and has the
\(\mathrm{DRat}\) data with index zero deleted.
\begin{proof}
Put a nonzero line vector first in a basis.  The matrix of \(r\)
is upper block triangular, with top left entry \(\psi_0\),
zero lower left block, and quotient block \(r'\).
The latter entries are continuous matrix coordinate functions;
its inverse is the same block of \(r(g^{-1})\), so it gives
a continuous quotient representation.

Fix any original bridge and enlarge coefficients to contain the
basis, the character, and the representation models.  Exactness
of de Rham comparison makes the quotient de Rham, and strictness
of the filtered exact sequence makes every labeled jump
multiplicity additive.  Cancellation of the singleton
\(\{h_0(\tau)\}\) gives the required quotient multiset,
even when weights repeat.  Choose a finite extension where all
three representations are semistable, using Berger for the de
Rham representations.  The exact WD sequence is upper block
triangular, hence its characteristic polynomial is the product
of those of the line and quotient.  Cancel the nonzero monic
polynomial \(X-\alpha_0(s)\) in \(\mathbb C[X]\).
This gives the required quotient polynomial for every Weil
element.  The comparison and finite descent laws identify this
calculation with the original bridge's parameter.  Thus every
clause of \(\mathrm{DRat}\) descends, not only its multiset
of Frobenius eigenvalues.
\end{proof}
% END OFL12

% BEGIN OFL13
\paragraph{OFL13. The first ordinary eigenline.}
\label{claim:OFL13}
Let \(n\ge1\), and suppose the ordinary characters
\((\omega_j,h_j,\psi_j)\) satisfy
\(h_0(\tau)<\cdots<h_{n-1}(\tau)\) for every embedding
\(\tau\), with pairwise distinct inertia restrictions of
\(\omega_j\circ a\).  If \(r\) has \(\mathrm{DRat}\)
weights \(\{h_j(\tau)\}_j\) and polynomials
\(\prod_j(X-\omega_j(a(s)))\), then \(r\) has a nonzero
eigenvector with character \(\psi_0\).
\begin{proof}
Twist \(r\) by \(\psi_0^{-1}\).  OFL1 and OFL4 give
the exact rank-one inverse data, and OFL11 gives the twisted
weights \(h_j(\tau)-h_0(\tau)\) and Weil characters
\(\mu_j=(\omega_j\omega_0^{-1})\circ a\).
They are pairwise distinct on inertia, because multiplying by the
same nonzero character preserves inequalities.  The zeroth
character is trivial.  Every shifted weight is nonnegative by
the stated order.  OFL10 gives a nonzero invariant vector \(v\)
of the twist.  The equality
\(\psi_0(g)^{-1}r(g)v=v\) is precisely
\(r(g)v=\psi_0(g)v\) for every \(g\in D\).
This constructs the first line in its prescribed order; it does
not choose an arbitrary inertia eigenline.
\end{proof}
% END OFL13

% BEGIN OFL14
\paragraph{OFL14. The full ordered invariant flag.}
\label{claim:OFL14}
Under the hypotheses of OFL13, allowing also \(n=0\), there is
a complete \(D\)-stable flag
\[
 0=V_0\subset V_1\subset\cdots\subset V_n
             =\overline{\mathbb Q}_\ell^n,
 \qquad D\text{ acts on }V_{j+1}/V_j\text{ by }\psi_j.
 \tag{OFL.11}
\]
\begin{proof}
Induct on \(n\).  In rank zero the empty flag has the required
property.  For positive rank, OFL13 gives the first eigenline
\(V_1\).  Complete its nonzero vector to a basis, so that the
quotient is the continuous representation in OFL12.  OFL4 provides
the rank-one \(\mathrm{DRat}\) for \(\psi_0\), and OFL12
gives the quotient data with index zero removed.  The remaining
weights are still strictly increasing, and their smooth inertia
characters are still pairwise distinct.  Induction gives their
ordered flag in the quotient.  Taking its inverse images and
preceding them by \(0\subset V_1\) gives (OFL.11).

The concatenated basis consists of the first eigenvector followed
by lifts of the quotient basis.  In it every matrix is upper
triangular with diagonal
\((\psi_0(g),\ldots,\psi_{n-1}(g))\).  Its change-of-basis
matrix is one invertible matrix over the original coefficient
closure, independent of \(g\).  This is exactly the source
\(\mathrm{HasFlag}(r,\psi)\), with the original order and
without a completion of the coefficient field.
\end{proof}
% END OFL14

% BEGIN OFL15
\paragraph{OFL15. Distinct unit characters give the original inertia hypothesis.}
\label{claim:OFL15}
The original pairwise distinct restrictions of \(\omega_j\) to
\(\mathcal O_{F_w}^\times\) imply the inertia distinctness used
in OFL5--OFL14.
\begin{proof}
For \(i\ne j\), choose a unit \(u\) with
\(\omega_i(u)\ne\omega_j(u)\).  Surjectivity of the
original local Artin map gives \(s\in W\) with \(a(s)=u\).
Its degree is zero, because
\(0=\operatorname{ord}_w(u)=-\deg(s)\).  Thus \(s\)
is the image of an inertia element, which supplies the required
inequality.  This argument uses the exact source reciprocity
sign and does not require an additional choice of uniformizer.
\end{proof}
% END OFL15

% BEGIN OFL16
\paragraph{OFL16. The literal \texttt{ordFlag\_alt} endpoint.}
\label{claim:OFL16}
The original \(\mathrm{LiteratureRankNHypothesis}\) implies
the complete original \(\mathrm{OrdFlag}\), at every allowed
prime and rank, on every original field, place, Artin map, and
coefficient isomorphism.
\begin{proof}
Unpack an arbitrary test datum in \(\mathrm{OrdFlag}\).
Its characters satisfy (OFL.3), its labeled weights are strictly
increasing in the stated order, its unit characters are distinct,
and its continuous representation satisfies the exact
\(\mathrm{DRat}\) polynomial identity.  OFL15 converts the
unit condition to inertia distinctness.  OFL14 gives
\(\mathrm{HasFlag}\) with exactly the supplied characters.
For every inducing bridge, OFL10 applied to the original,
untwisted representation gives potential crystallinity.  These
are the two conclusions of the original \(\mathrm{OrdFlag}\).
The original odd-prime hypothesis, when present in the consumer,
is retained; none of the steps discards the prime \(3\) or
requires a new lower bound on it.  No semisimplicity hypothesis
on \(r\) is used or added.
\end{proof}
% END OFL16

\paragraph{Boundary of this chapter.}
The original Hecke eigencharacter calculation is still needed to
produce the \(\mathrm{OrdChar}\) inputs at arithmetic Hecke
points.  That calculation is supplied separately by the HEP
chapter.  This chapter proves their complete ordered-flag consumer,
while retaining the named published and period-ring foundations
listed at the start.
